\section{Ownership cancellation and failure propagation}\label{sec:C21:ownership}
Every delegated operation needs an owner responsible for its budget, authority, cancellation, and final integration. A manager that starts workers and then terminates can leave active processes with no accountable controller. Ownership should be represented in execution state, with child operations linked to parent task and policy epochs. A child cannot acquire permissions merely because its parent was authorized to delegate reasoning.

Delegation authority is narrower than tool authority. A worker may be allowed to inspect a repository but not publish a branch, or to retrieve enterprise documents but not export them. Passing an unrestricted service credential to every worker defeats role isolation. Capability-scoped credentials and application-enforced tool registries make the distinction executable. The language model's stated role is useful task guidance, but it is not an access-control boundary.

Cancellation has two meanings: stop producing additional work and stop all ongoing effects. The former may be achievable immediately; the latter can require tool-specific termination or reconciliation. If one worker completes the task, cancelling redundant workers should not discard evidence of requests already dispatched. The system needs an outstanding-operation set and a rule for handling late responses. Otherwise a result can arrive after the manager has finalized and silently mutate shared memory.

Failure propagation should distinguish dependent and independent branches. If a shared prerequisite is invalidated, every descendant that used it may require reevaluation. An unrelated worker can continue. This can be implemented with artifact dependency digests rather than broad restart signals. A global restart is simpler but increases cost and can repeat external effects. Overly narrow invalidation is cheaper but risks integrating work derived from obsolete premises.

The tail behavior of a team often depends on stragglers. Waiting for every worker maximizes coverage but can make latency dominated by one failed dependency. Quorum completion can reduce latency but changes the evidence population and may systematically omit difficult or dissenting analyses. The stopping rule should specify whether missing workers count as failures, censored observations, or intentionally unused candidates. A post hoc choice of whichever stopping rule makes the run look best is evaluation leakage.

Fault-injection tests should combine worker loss, delayed messages, policy revocation, and shared-store outages. Independent single faults are easier to diagnose, but correlated failures are operationally important: one tool outage can affect all workers, while one poisoned shared memory can redirect the whole team. Report containment and benign completion together. A coordinator that avoids every incident by refusing delegation has not demonstrated useful orchestration.

